% Manual de usuario - SIIIS V3
% Compilar con pdflatex (dos pasadas). Las capturas van en img/; si falta alguna, se muestra "Captura pendiente".
\documentclass[11pt,a4paper]{article}
\IfFileExists{siiis.sty}{\usepackage{siiis}}{\usepackage{../siiis}}
\renewcommand{\siiiscarpeta}{manual-usuario}
\graphicspath{{img/}{manual-usuario/img/}}

\siiisdatos{Manual de usuario}{SIIIS-MU-001}{1.0}{28 de noviembre de 2026}

% Nombre de un boton, pestana o enlace tal como aparece en pantalla
\newcommand{\boton}[1]{\textbf{\textcolor{siiisazul}{#1}}}

\begin{document}

\siiisportada{Manual de usuario}{Sitio web SIIIS V3: gu\'ia para visitantes, integrantes y administradores}

\siiiscontrolversiones{%
  1.0 & 28/11/2026 & Versi\'on inicial del manual de usuario. & Equipo de Desarrollo SIIIS\\
}

\tableofcontents
\clearpage

% =================================================================================================================
\section{Introducci\'on}
% =================================================================================================================

Este manual explica c\'omo usar el sitio web del Semillero de Investigaci\'on SIIIS, disponible en
\url{https://siiis.com.co}. Est\'a escrito para tres tipos de personas:

\begin{itemize}
  \item \textbf{Visitantes}: cualquier persona que entra al sitio sin iniciar sesi\'on.
  \item \textbf{Integrantes}: estudiantes y docentes del semillero que tienen una cuenta.
  \item \textbf{Administradores}: quienes gestionan los usuarios, el contenido y las solicitudes desde el panel de
    administraci\'on.
\end{itemize}

El sitio funciona en computadores, tabletas y celulares. En el celular, el men\'u principal se abre con el bot\'on de
las tres l\'ineas, en la parte superior.

\captura[0.95\linewidth]{inicio.png}{P\'agina de inicio en un computador.}
\captura[0.4\linewidth]{inicio-celular.png}{P\'agina de inicio en un celular, con el men\'u abierto.}

% =================================================================================================================
\section{Roles y qu\'e puede hacer cada uno}
% =================================================================================================================

Cada cuenta tiene un rol. El rol lo asigna un administrador; nadie puede cambiar su propio rol.

{\small
\begin{longtable}{@{}p{3.6cm}p{10.4cm}@{}}
\toprule
\textbf{Rol} & \textbf{Qu\'e puede hacer}\\
\midrule
\endhead
\textbf{Visitante} (sin cuenta) & Navegar el sitio; ver integrantes, proyectos, art\'iculos y noticias; enviar la
  solicitud de inscripci\'on; escribir al semillero desde Contacto; dejar una rese\~na (se publica cuando la revisan);
  ir al blog.\\
\textbf{Estudiante} (\cod{estudiante}) & Todo lo del visitante, m\'as: iniciar sesi\'on; editar su perfil, cambiar o
  quitar su foto y cambiar su contrase\~na; crear proyectos, art\'iculos y noticias (quedan en revisi\'on) y editarlos o
  eliminarlos mientras siguen en revisi\'on; ver sus borradores y aquellos en los que aparece como integrante; dejar
  rese\~nas que se publican de inmediato con su nombre.\\
\textbf{Estudiante l\'ider} (\cod{estudiante_lider}) & Lo mismo que un estudiante. Hoy el sitio no le da permisos
  adicionales.\\
\textbf{Docente director} (\cod{docente_director}) & Todo lo anterior, m\'as el panel de administraci\'on con dos
  secciones: \boton{Contenido en revisi\'on} (publicar, ocultar, editar y eliminar cualquier proyecto, art\'iculo o
  noticia) y \boton{Rese\~nas} (publicar, ocultar o borrar rese\~nas).\\
\textbf{Administrador} (\cod{administrador}) & Todo lo del docente director, m\'as: \boton{Usuarios} (crear, editar,
  desactivar y reactivar cuentas de docentes directores, estudiantes l\'ideres y estudiantes), \boton{Inscripciones},
  \boton{Mensajes de contacto} e \boton{Inicio} (fotos del carrusel y video de presentaci\'on).\\
\textbf{Superadministrador} (\cod{superadministrador}) & Todo lo del administrador, y adem\'as gestiona cuentas de
  administradores y de otros superadministradores. El sitio no deja desactivar ni quitar el rol al \'ultimo
  superadministrador activo.\\
\bottomrule
\end{longtable}
}

% =================================================================================================================
\section{Como visitante}
% =================================================================================================================

\subsection{Navegar el sitio}
El men\'u principal tiene: \boton{Inicio}, \boton{Acerca de}, \boton{Integrantes}, \boton{Proyectos},
\boton{Art\'iculos}, \boton{Noticias}, \boton{Blog}, \boton{Contacto} e \boton{Inscripci\'on}. En la parte inferior de
cada p\'agina est\'an los enlaces a la pol\'itica de privacidad y a los t\'erminos de uso.

\begin{itemize}
  \item \textbf{Inicio}: carrusel de fotos del semillero, video de presentaci\'on y noticias recientes.
  \item \textbf{Acerca de}: qui\'enes somos, la direcci\'on del semillero y los equipos de las versiones anteriores del
    sitio.
\end{itemize}

Al cambiar de p\'agina aparece por un momento el logo del semillero girando como una moneda: es la animaci\'on de carga.

\subsection{Ver a los integrantes}
\begin{enumerate}
  \item Pulse \boton{Integrantes} en el men\'u.
  \item Recorra el listado. Puede filtrar por rol con la lista desplegable.
  \item Pulse el nombre de una persona para ver su perfil: programa acad\'emico, semestre, biograf\'ia y sus
    publicaciones y proyectos.
  \item A un lado del perfil (debajo, en el celular) est\'a la columna \boton{Art\'iculos de inter\'es}, con los
    \'ultimos art\'iculos de otros integrantes.
\end{enumerate}
El correo y el tel\'efono de los integrantes nunca se muestran en el sitio.

\captura{integrantes.png}{Listado de integrantes.}
\captura{perfil-integrante.png}{Perfil de un integrante con sus publicaciones y la columna ``Art\'iculos de inter\'es''.}

\subsection{Ver proyectos, art\'iculos y noticias}
\begin{enumerate}
  \item Pulse \boton{Proyectos}, \boton{Art\'iculos} o \boton{Noticias} en el men\'u.
  \item Use los filtros de la parte superior si los hay. En proyectos puede filtrar por estado (\emph{En ejecuci\'on} o
    \emph{Finalizado}).
  \item Pulse una tarjeta para abrir la publicaci\'on completa. All\'i ver\'a el texto, el documento PDF (si lo tiene)
    con el bot\'on \boton{Descargar PDF}, las im\'agenes y los integrantes que participaron.
\end{enumerate}

\captura{proyectos.png}{Listado de proyectos.}
\captura{detalle-proyecto.png}{Detalle de un proyecto con su texto y el visor del PDF.}

\subsection{Enviar la solicitud de inscripci\'on}
\begin{enumerate}
  \item Pulse \boton{Inscripci\'on} en el men\'u.
  \item Complete el formulario:
    \begin{itemize}
      \item \textbf{Nombre completo}.
      \item \textbf{Correo electr\'onico}: debe ser su correo institucional \cod{@uptc.edu.co}.
      \item \textbf{C\'odigo estudiantil} (opcional).
      \item \textbf{Programa acad\'emico} y \textbf{Semestre}.
      \item \textbf{\'Areas de inter\'es}: marque una o varias.
      \item \textbf{Carta de intenci\'on o motivaci\'on}: cu\'entenos por qu\'e quiere unirse.
      \item \textbf{Hoja de vida (PDF)}, opcional: pulse \boton{Seleccionar archivo} y elija un PDF de hasta 3 MB.
        Debajo ver\'a el nombre del archivo elegido.
      \item Marque la casilla de disponibilidad para las reuniones semanales (obligatoria).
    \end{itemize}
  \item Pulse \boton{Enviar solicitud}. Ver\'a un mensaje de confirmaci\'on y le escribiremos al correo que indic\'o con
    la respuesta.
\end{enumerate}
Solo puede haber una solicitud pendiente por correo. Si su solicitud ya fue aceptada o rechazada, puede enviar otra.

\captura{inscripcion.png}{Formulario de inscripci\'on.}

\subsection{Escribir al semillero (Contacto)}
\begin{enumerate}
  \item Pulse \boton{Contacto} en el men\'u.
  \item Escriba su \textbf{Nombre}, \textbf{Correo electr\'onico}, \textbf{Asunto} y \textbf{Mensaje}.
  \item Pulse \boton{Enviar}. El equipo del semillero le responder\'a al correo que indic\'o.
\end{enumerate}

\captura{contacto.png}{P\'agina de contacto.}

\subsection{Dejar una rese\~na}
\begin{enumerate}
  \item Baje hasta la secci\'on de rese\~nas de la p\'agina de inicio (o abra la p\'agina de rese\~nas desde el bot\'on
    de esa secci\'on).
  \item En \boton{Deja tu rese\~na}, escriba \textbf{Tu nombre}, elija una \textbf{Calificaci\'on} (opcional) y escriba
    el \textbf{Comentario}.
  \item Pulse \boton{Enviar rese\~na}.
\end{enumerate}
Las rese\~nas de visitantes se publican cuando el equipo las revisa. Si usted es integrante y tiene la sesi\'on
iniciada, su rese\~na se publica de inmediato con su nombre.

\captura{resenas.png}{Rese\~nas del semillero y formulario para dejar una.}

\subsection{Ir al blog}
La pesta\~na \boton{Blog} del men\'u abre en otra pesta\~na el blog del semillero (\url{https://semillerosiiis.blogspot.com/}).

% =================================================================================================================
\section{Como integrante}
% =================================================================================================================

Para usar estas funciones necesita una cuenta. Las cuentas las crea un administrador (normalmente despu\'es de aceptar
su solicitud de inscripci\'on); no hay registro libre en el sitio.

\subsection{Iniciar sesi\'on}
\textbf{Con contrase\~na:}
\begin{enumerate}
  \item Pulse \boton{Iniciar sesi\'on}, arriba a la derecha.
  \item Escriba su \textbf{Correo electr\'onico} y su \textbf{Contrase\~na}.
  \item Si est\'a en su propio equipo, puede marcar \boton{Recordarme} para que la sesi\'on siga abierta aunque cierre
    el navegador (hasta un m\'aximo de 8 horas).
  \item Pulse \boton{Ingresar}.
\end{enumerate}

\textbf{Con Google:}
\begin{enumerate}
  \item Pulse \boton{Iniciar sesi\'on} y luego el bot\'on \boton{Iniciar sesi\'on con Google}.
  \item Elija su cuenta institucional \cod{@uptc.edu.co}.
\end{enumerate}
Con Google solo pueden entrar cuentas \cod{@uptc.edu.co} que ya tengan usuario en el semillero. Si su cuenta no est\'a
registrada, el sitio se lo indica y le ofrece ir a la inscripci\'on.

Con la sesi\'on iniciada, al pulsar su nombre arriba a la derecha ver\'a: \boton{Mi perfil}, \boton{Mis contenidos},
\boton{Administraci\'on} (solo si su rol tiene panel) y \boton{Cerrar sesi\'on}.

\captura[0.6\linewidth]{login.png}{Ventana de inicio de sesi\'on con correo, contrase\~na y el bot\'on de Google.}

\subsection{Recuperar la contrase\~na}
\begin{enumerate}
  \item En la ventana de inicio de sesi\'on, pulse \boton{¿Olvidaste tu contrase\~na?}.
  \item Escriba su correo y pulse \boton{Enviar enlace}.
  \item Revise su correo (tambi\'en la carpeta de correo no deseado). El enlace sirve una sola vez y vence en 30 minutos.
  \item Abra el enlace. En la p\'agina \boton{Restablecer contrase\~na}, escriba la \textbf{Nueva contrase\~na}: al menos
    8 caracteres, con al menos una letra y un n\'umero, y distinta de la actual.
  \item Pulse \boton{Guardar contrase\~na} e inicie sesi\'on con la nueva.
\end{enumerate}
Por seguridad, el sitio muestra el mismo mensaje exista o no el correo, y al cambiar la contrase\~na se cierran las
sesiones abiertas en otros equipos.

\captura[0.6\linewidth]{recuperar-contrasena.png}{Solicitud del enlace para recuperar la contrase\~na.}
\captura[0.6\linewidth]{restablecer.png}{P\'agina para escribir la nueva contrase\~na.}

\subsection{Cierre autom\'atico a las 8 horas}
Por seguridad, toda sesi\'on dura como m\'aximo 8 horas. Pasado ese tiempo el sitio la cierra solo, aunque est\'e en
otra pesta\~na, y lo lleva a iniciar sesi\'on con el aviso ``Tu sesi\'on expir\'o. Inicia sesi\'on de nuevo''. Tambi\'en
se cierra si un administrador cambia su rol o desactiva su cuenta. Lo que no haya guardado en ese momento se pierde, as\'i
que conviene guardar con frecuencia.

\captura[0.6\linewidth]{sesion-expirada.png}{Aviso de sesi\'on expirada.}

\subsection{Editar el perfil y la foto}
\begin{enumerate}
  \item Pulse su nombre y luego \boton{Mi perfil}.
  \item En \boton{Mis datos}, actualice su nombre, apellido, programa acad\'emico, semestre, tel\'efono (no es p\'ublico)
    y biograf\'ia. Pulse \boton{Guardar cambios}.
  \item Para \textbf{subir o cambiar la foto}, pulse \boton{Subir foto} o \boton{Cambiar foto} y elija una imagen JPG,
    PNG o WebP de hasta 2 MB. Se recorta autom\'aticamente en un cuadrado.
  \item Para \textbf{quitar la foto}, pulse \boton{Quitar foto} y confirme en el di\'alogo. En su lugar se mostrar\'an
    sus iniciales, en el perfil y en el men\'u.
  \item Para \textbf{cambiar la contrase\~na}, use la secci\'on \boton{Cambiar contrase\~na}: escriba la actual, la
    nueva y su confirmaci\'on.
\end{enumerate}

\captura{mi-perfil.png}{Pantalla ``Mi perfil'' con la foto, los botones ``Cambiar foto'' y ``Quitar foto'' y los datos.}
\captura[0.6\linewidth]{quitar-foto.png}{Di\'alogo de confirmaci\'on para quitar la foto.}

\subsection{Crear, editar y eliminar proyectos, art\'iculos y noticias}

\subsubsection*{Crear}
\begin{enumerate}
  \item Pulse su nombre y luego \boton{Mis contenidos}.
  \item Pulse \boton{Nuevo proyecto}, \boton{Nuevo art\'iculo} o \boton{Nueva noticia}.
  \item Complete el \textbf{T\'itulo} y el \textbf{Resumen} (aparece en los listados).
  \item Agregue el contenido. Puede usar \textbf{texto}, un \textbf{PDF} o \textbf{ambos}:
    \begin{itemize}
      \item \textbf{Texto}: escr\'ibalo en el editor de \textbf{Contenido} (ver la gu\'ia m\'as abajo).
      \item \textbf{PDF}: pulse \boton{Seleccionar PDF} y elija un PDF de hasta 3 MB.
    \end{itemize}
  \item Si otras personas participaron, b\'usquelas en \textbf{Integrantes} y agr\'eguelas (hasta 20).
  \item Pulse \boton{Crear}.
\end{enumerate}
Lo que crea un integrante queda \textbf{En revisi\'on} y oculto al p\'ublico hasta que un revisor (docente director,
administrador o superadministrador) lo publique. Una vez creado, puede agregar im\'agenes con \boton{Subir imagen} (JPG,
PNG o WebP de hasta 5 MB; m\'aximo 10 archivos por publicaci\'on).

\subsubsection*{Editar y eliminar}
\begin{enumerate}
  \item En \boton{Mis contenidos}, busque la publicaci\'on y pulse \boton{Editar}.
  \item Haga los cambios y pulse \boton{Guardar cambios}. Para cambiar el PDF elija otro con \boton{Reemplazar PDF};
    para quitarlo, pulse \boton{Quitar PDF} (solo si la publicaci\'on conserva texto).
  \item Para eliminarla, abra la publicaci\'on y pulse \boton{Eliminar}; confirme en el di\'alogo.
\end{enumerate}
El autor puede editar y eliminar su publicaci\'on solo mientras sigue \textbf{En revisi\'on}. Despu\'es de publicada,
los cambios los hace un revisor. Las personas agregadas como integrantes ven la publicaci\'on, pero no la editan.

\subsubsection*{Gu\'ia r\'apida del editor}
El editor tiene dos pesta\~nas: \boton{Escribir} y \boton{Vista previa}. En \boton{Escribir}, la barra de botones da
formato al texto seleccionado:

\begin{center}
\begin{tabular}{@{}p{3.6cm}p{10.4cm}@{}}
  \toprule
  \textbf{Bot\'on} & \textbf{Qu\'e hace}\\
  \midrule
  Negrita & Resalta el texto en \textbf{negrita}.\\
  Cursiva & Pone el texto en \emph{cursiva}.\\
  T\'itulo y Subt\'itulo & Convierten la l\'inea en un t\'itulo o un subt\'itulo.\\
  Lista con vi\~netas & Convierte las l\'ineas en una lista con puntos.\\
  Lista numerada & Convierte las l\'ineas en una lista 1, 2, 3\ldots\\
  Enlace & Inserta un enlace: escriba el texto y reemplace la direcci\'on que aparece entre par\'entesis.\\
  Cita & Convierte la l\'inea en una cita.\\
  \bottomrule
\end{tabular}
\end{center}

La pesta\~na \boton{Vista previa} muestra c\'omo se ver\'a el texto publicado, sin guardarlo.

\captura{mis-contenidos.png}{Pantalla ``Mis contenidos'' con los botones para crear y la lista de publicaciones propias.}
\captura{editor-documento.png}{Editor de una publicaci\'on: t\'itulo, resumen, contenido, PDF e integrantes.}
\captura{editor-vista-previa.png}{Pesta\~na ``Vista previa'' del editor.}

\subsection{Art\'iculos de inter\'es y vista previa}
\begin{enumerate}
  \item Abra el perfil de un integrante desde \boton{Integrantes}.
  \item Pulse cualquier publicaci\'on del perfil o de la columna \boton{Art\'iculos de inter\'es}.
  \item Se abre una ventana con la vista previa: el texto y el visor del PDF, sin salir del perfil.
  \item Pulse \boton{Abrir p\'agina completa} para ir a la publicaci\'on, o \boton{Cerrar} (o la tecla Esc) para volver.
\end{enumerate}

\captura{vista-previa-publicacion.png}{Vista previa de una publicaci\'on desde el perfil de un integrante.}

% =================================================================================================================
\section{Como administrador y superadministrador}
% =================================================================================================================

\subsection{Entrar al panel}
Pulse su nombre y luego \boton{Administraci\'on}. El panel muestra solo las secciones que su rol permite:

\begin{center}
\begin{tabular}{@{}p{4.4cm}p{9.6cm}@{}}
  \toprule
  \textbf{Secci\'on} & \textbf{Para qu\'e sirve}\\
  \midrule
  Contenido en revisi\'on & Publicar, ocultar o borrar proyectos, art\'iculos y noticias.\\
  Rese\~nas & Aprobar u ocultar rese\~nas y borrar las ofensivas.\\
  Inicio & Fotos del carrusel y video de presentaci\'on de la p\'agina de inicio.\\
  Usuarios & Crear cuentas, cambiar roles, desactivar y reactivar integrantes.\\
  Inscripciones & Revisar solicitudes y sus documentos; aceptar o rechazar.\\
  Mensajes de contacto & Bandeja de mensajes con le\'idos y no le\'idos.\\
  \bottomrule
\end{tabular}
\end{center}
El docente director ve solo \boton{Contenido en revisi\'on} y \boton{Rese\~nas}.

\captura{panel.png}{Panel de administraci\'on de un administrador.}

\subsection{Contenido en revisi\'on}
\begin{enumerate}
  \item Entre a \boton{Contenido en revisi\'on}. Elija \boton{Proyectos} o \boton{Art\'iculos y noticias} y, si quiere,
    filtre por visibilidad.
  \item Abra una publicaci\'on. Seg\'un su estado ver\'a los botones:
    \begin{itemize}
      \item \boton{Publicar}: la hace visible al p\'ublico.
      \item \boton{Ocultar} y \boton{Mostrar}: la quitan o la devuelven al sitio p\'ublico sin borrarla.
      \item \boton{Marcar finalizado}: para proyectos en ejecuci\'on que ya terminaron.
      \item \boton{Devolver a revisi\'on}: la oculta y permite que el autor la vuelva a editar.
      \item \boton{Editar} y \boton{Eliminar} (con confirmaci\'on; no se puede deshacer).
    \end{itemize}
\end{enumerate}
Una publicaci\'on necesita texto, PDF o ambos para poder publicarse.

\captura{admin-contenido.png}{Listado de ``Contenido en revisi\'on''.}

\subsection{Moderar rese\~nas}
\begin{enumerate}
  \item Entre a \boton{Rese\~nas}. Las de visitantes llegan ocultas.
  \item Pulse \boton{Publicar} para mostrarla u \boton{Ocultar} para retirarla del sitio.
  \item Use \boton{Borrar} solo para spam o contenido ofensivo: el navegador pide confirmar y no se puede deshacer.
\end{enumerate}

\captura{admin-resenas.png}{Moderaci\'on de rese\~nas.}

\subsection{Usuarios}
\textbf{Crear una cuenta:}
\begin{enumerate}
  \item Entre a \boton{Usuarios} y pulse \boton{Crear usuario}.
  \item Escriba el \textbf{Correo}, una \textbf{Contrase\~na inicial} (8 o m\'as caracteres con letra y n\'umero), el
    \textbf{Rol} y los datos del perfil.
  \item Pulse \boton{Crear usuario}. La cuenta queda activa; entregue la contrase\~na inicial a la persona, que puede
    cambiarla en \boton{Mi perfil}.
\end{enumerate}

\textbf{Editar:} en la fila de la persona, pulse \boton{Editar}, cambie sus datos o su rol y pulse
\boton{Guardar cambios}. Su propia cuenta se edita desde \boton{Mi perfil}.

\textbf{Desactivar y reactivar:} en el listado, pulse \boton{Desactivar} y confirme. La persona deja de poder entrar y
de aparecer en el sitio, y sus sesiones se cierran; sus datos no se borran. Para devolverle el acceso, pulse
\boton{Reactivar}.

Un administrador gestiona docentes directores, estudiantes l\'ideres y estudiantes; el superadministrador gestiona
todos los roles. Nadie puede cambiar su propio rol ni desactivarse a s\'i mismo.

\captura{admin-usuarios.png}{Listado de usuarios con los botones ``Crear usuario'', ``Desactivar'' y ``Reactivar''.}
\captura{admin-usuario-nuevo.png}{Formulario para crear un usuario.}

\subsection{Inscripciones}
\begin{enumerate}
  \item Entre a \boton{Inscripciones}. Filtre por \boton{Pendientes}, \boton{Aceptadas}, \boton{Rechazadas} o
    \boton{Todas}.
  \item Abra una solicitud para ver los datos del aspirante, sus \'areas de inter\'es y la carta de motivaci\'on.
  \item En \boton{Documentos (hoja de vida)} est\'an los archivos adjuntos. Los enlaces son privados y vencen en 1 hora;
    si vencieron, pulse \boton{Renovar enlaces}.
  \item Pulse \boton{Aceptar} o \boton{Rechazar}. Al aspirante le llega un correo con la respuesta. Una solicitud se
    resuelve una sola vez.
  \item Aceptar \textbf{no} crea la cuenta: despu\'es pulse \boton{Crear cuenta con estos datos}, revise los datos que
    se copian de la solicitud, defina una contrase\~na inicial y env\'iesela al aspirante.
\end{enumerate}
Las solicitudes rechazadas se conservan seis meses y luego las borra el equipo t\'ecnico.

\captura{admin-inscripciones.png}{Listado de inscripciones.}
\captura{admin-inscripcion-detalle.png}{Detalle de una solicitud con la hoja de vida y los botones ``Aceptar'' y ``Rechazar''.}

\subsection{Mensajes de contacto}
\begin{enumerate}
  \item Entre a \boton{Mensajes de contacto}. Filtre por \boton{No le\'idos}, \boton{Le\'idos} o \boton{Todos}.
  \item Al abrir un mensaje, queda marcado como le\'ido.
  \item Pulse \boton{Responder} para escribirle al remitente desde su programa de correo.
\end{enumerate}

\captura{admin-mensajes.png}{Bandeja de mensajes de contacto.}

\subsection{Multimedia del inicio}
Entre a \boton{Inicio} en el panel. Los cambios se ven en la p\'agina de inicio en unos minutos.

\textbf{Carrusel de fotos:}
\begin{enumerate}
  \item Para subir una foto: pulse \boton{Seleccionar foto}, elija una imagen JPG, PNG o WebP de hasta 5 MB, escriba
    una \textbf{Descripci\'on} (la leen los lectores de pantalla), indique la \textbf{Posici\'on} si quiere (opcional)
    y pulse \boton{Subir foto}.
  \item Para cambiar el orden, use \boton{Antes} y \boton{Despu\'es} en cada foto.
  \item Para cambiar la descripci\'on, ed\'itela y gu\'ardela.
  \item Para quitar una foto, pulse \boton{Borrar} y confirme. Se elimina tambi\'en el archivo y no se puede deshacer.
\end{enumerate}

\textbf{Video de presentaci\'on:} solo puede haber uno.
\begin{enumerate}
  \item Pulse \boton{Seleccionar video} y elija un video MP4, WebM o MOV de hasta 50 MB.
  \item Pulse \boton{Reemplazar video} (o \boton{Subir video} si todav\'ia no hay ninguno). El video anterior se borra
    solo si el nuevo se sube bien. Un video grande puede tardar varios minutos en subir: no cierre la p\'agina.
\end{enumerate}

\captura{admin-inicio.png}{Administraci\'on del carrusel y del video de presentaci\'on.}

% =================================================================================================================
\section{Preguntas frecuentes}
% =================================================================================================================

\begin{description}[style=nextline]
  \item[¿C\'omo obtengo una cuenta?] Env\'ie la solicitud desde \boton{Inscripci\'on}. Si la aceptan, un administrador
    crea su cuenta y le entrega una contrase\~na inicial.
  \item[Google dice que mi cuenta no est\'a registrada.] Solo entran con Google las cuentas \cod{@uptc.edu.co} que ya
    tienen usuario en el semillero. Pida a un administrador que revise su cuenta o env\'ie la inscripci\'on.
  \item[No me lleg\'o el correo para recuperar la contrase\~na.] Revise la carpeta de correo no deseado. El enlace vence
    en 30 minutos y solo se env\'ia a cuentas activas. Si sigue sin llegar, escriba a soporte.
  \item[El sitio dice ``Demasiados intentos''.] Por seguridad hay un l\'imite de intentos. Espere unos minutos y vuelva a
    intentarlo.
  \item[Mi publicaci\'on no aparece en el sitio.] Lo que crea un integrante queda \emph{En revisi\'on} hasta que un
    revisor lo publique.
  \item[Ya no puedo editar mi publicaci\'on.] El autor solo edita mientras est\'a en revisi\'on. Pida a un revisor que
    haga el cambio o que la devuelva a revisi\'on.
  \item[Me sac\'o de la sesi\'on sin avisar.] La sesi\'on dura como m\'aximo 8 horas. Tambi\'en se cierra si cambian su
    rol o su cuenta se desactiva.
  \item[No se sube mi archivo.] Revise el formato y el tama\~no: fotos de perfil hasta 2 MB; PDF hasta 3 MB; im\'agenes
    hasta 5 MB; video del inicio hasta 50 MB.
  \item[Hice un cambio y no lo veo.] Algunas p\'aginas tardan unos minutos en actualizarse. Recargue la p\'agina; si
    persiste, use la recarga forzada (Ctrl + F5).
  \item[Quiero que eliminen mi cuenta.] P\'idalo a un administrador: la cuenta se desactiva y deja de aparecer en el
    sitio.
\end{description}

% =================================================================================================================
\section{Contacto de soporte}
% =================================================================================================================

Para problemas con el sitio, escriba a \href{mailto:siiisweb@gmail.com}{\textbf{siiisweb@gmail.com}} indicando qu\'e
estaba haciendo, en qu\'e p\'agina y, si puede, una captura de pantalla del mensaje que vio. Nunca env\'ie su
contrase\~na: el equipo del semillero no la pedir\'a.

\subsection*{Equipo de desarrollo}
\siiisequipo

\end{document}
